\section{Related Work}
\label{sec:related}

\subsection{Image scene text editing}
\label{sec:related:ste}

Scene text editing (STE)---modifying the textual content of a
natural image while preserving background, typography, and
lighting---has grown into a small but active sub-field of
image-to-image translation. Early learned approaches such as
SRNet~\citep{wu2019srnet} cast STE as the composition of an
inpainting module, a text-style transfer module, and a fusion
module, trained adversarially with synthetic source/target pairs.
MOSTEL~\citep{qu2023mostel} extends this line of work with
stroke-level supervision for finer glyph control. More recent
diffusion-based systems have come to dominate the area:
AnyText~\citep{tuo2024anytext} and its successor
AnyText2~\citep{tuo2025anytext2} combine a Stable Diffusion
backbone~\citep{rombach2022stablediffusion} with a text-writing
network that conditions generation on a target string, producing
high-quality multilingual scene text. Our pipeline uses AnyText2 as
its editing backend (Section~\ref{sec:method:s3}). All of these
systems operate on single images and do not address the additional
challenges of temporal coherence, per-frame lighting and blur
adaptation, or propagation of a single edit across many frames of
video.

\subsection{Video scene text replacement}
\label{sec:related:video}

To our knowledge, the closest prior work to ours is
STRIVE~\citep{subramanian2021strive}, which proposes a
frontalize--edit--propagate--composite structure for replacing text
in video and introduces a Text Propagation Module consisting of a
Lighting Correction Module and a Blur Prediction Network. STRIVE's
code has, to our knowledge, never been released, which makes direct
reproduction impossible; our pipeline's overall shape is informed
by STRIVE's design philosophy, but the implementation is independent
and differs in important details (notably the use of a learned
alignment refiner described in Section~\ref{sec:method:refiner}).
Beyond STRIVE, we are not aware of a released end-to-end pipeline
for cross-language video text replacement, which complicates
quantitative comparison against prior work.

\subsection{Generative video editing}
\label{sec:related:generative}

A parallel and fast-growing paradigm is end-to-end generative video
editing. Commercial systems such as Runway's Gen-4
Aleph~\citep{runway2025aleph} and research systems such as
Sora~\citep{openai2024sora} rewrite arbitrary regions of video
conditioned on natural-language prompts, achieving striking
photorealism in lighting, texture, and motion. Reliable glyph-level
control---producing a visually correct rendering of a specific
target string, in the right script, grammar, and layout---remains an
open research problem for these systems at the time of writing. We
use Gen-4 Aleph as a qualitative comparison in
Section~\ref{sec:exp:qualitative}, and we emphasize that the
comparison is \emph{between paradigms}, not between competing
solutions to the same problem.

\subsection{Supporting vision components}
\label{sec:related:support}

Our pipeline relies on several established components. For scene
text detection and recognition we use PaddleOCR's
PP-OCRv4~\citep{paddleocr2024}, with a false-positive filter based
on wordfreq~\citep{wordfreq} that rejects OCR detections whose
recognized string has an anomalously low Zipf frequency. Temporal
propagation of text regions between OCR-detected frames uses
CoTracker3~\citep{karaev2024cotracker3}, a learned point tracker
that is substantially smoother and faster on scene text than
classical optical flow. For background isolation before lighting
correction we support two interchangeable inpainters: SRNet reused
as a learned text-removal network, and a segmentation-based pipeline
built on Hi-SAM~\citep{ye2024hisam}, which specializes the Segment
Anything Model~\citep{kirillov2023sam} to text strokes, followed by
Navier--Stokes inpainting. Compositing uses standard techniques:
feathered alpha blending, with a Poisson-blending
fallback~\citep{perez2003poisson} available for difficult seams.
The two networks we trained, described in
Sections~\ref{sec:method:bpn} and~\ref{sec:method:refiner}, are
built on a ResNet~\citep{he2016resnet} backbone and a simple
convolutional encoder respectively. Spatial-Temporal Transformer
Networks~\citep{zeng2020sttn} are a natural future direction for
replacing our classical homography with a learned temporal
frontalizer, but are not used in the current pipeline.
